\chapter{Introduction}

\section{Motivation}

{\small Robotic systems are increasingly moving out of fenced industrial cells and into shared, dynamic environments such as warehouses, hospitals, offices, and public spaces. According to the International Federation of Robotics, almost 200,000 professional service robots were sold in 2024, an increase of 9\% over the previous year~\cite{ifr2025servicerobots}. The largest share of these, 102,900 units, were built for transportation and logistics and consist mainly of mobile robots that move goods through their surroundings~\cite{ifr2025servicerobots}. At the same time, the global operational stock of industrial robots reached 4.66 million units in 2024~\cite{ifr2025robots}. Together, these figures show that robots are becoming a common part of environments that are also used by people.\par}

{\small Alongside this growth, robotic systems are increasingly shaped by artificial intelligence. The International Federation of Robotics identifies physical, analytic, and generative AI as key technological trends in robotics~\cite{ifr2025robots}. AI-based components allow robots to perceive their surroundings, interpret scenes, and adapt their behaviour. At the same time, they make robot behaviour harder to predict fully during design-time safety analysis~\cite{silvaneto2022aisafetyassurance}.\par}

{\small This development creates new challenges for safety and risk assessment. Conventional risk assessment is typically carried out before deployment and relies on assumptions about the system design, intended use, operating environment, and known hazards~\cite{iso12100}. These methods remain essential, but they reach their limits when mobile robots operate in changing environments, receive over-the-air software updates, or use adaptive AI-based components. In such cases, the situation the robot actually faces may differ from the one assumed during the original assessment. For this reason, automated and continuous risk assessment has been proposed, in which the assessment is updated while the system is running~\cite{grimmeisen2023continuousrisk}.\par}

{\small Such dynamic risk assessment needs a structured description of the current environment. It must know which objects are present, where they are, how they relate to each other, and which of them may be relevant to risk. A spatial risk-annotated scene graph is a suitable way to provide this information, as it combines perceived objects, their spatial context, and risk-related properties in a single machine-readable representation. In this thesis, the term \emph{risk-annotated scene graph} refers to a spatial scene graph whose object nodes are enriched with risk-relevant semantic and uncertainty information, including object labels, confidence values, mobility information, presence likelihood, and descriptive object properties. These annotations provide structured environmental context for downstream dynamic risk assessment rather than representing a complete risk assessment by themselves.\par}

\section{Relevance}

{\small Scene graphs have become an established way to represent robotic environments in a structured form. Hydra~\cite{hughes2022hydra} builds a layered 3D scene graph of objects, places, and rooms in real time and keeps it consistent through loop closure. However, Hydra does not identify objects itself. Its object layer depends on per-pixel semantic labels, which come from ground-truth annotations in simulation or from a segmentation network trained on a fixed set of classes~\cite{hughes2022hydra}. ConceptGraphs~\cite{gu2024conceptgraphs} removes this restriction by identifying objects with open-vocabulary foundation models, but it applies these models to every observation and represents objects through free-form descriptions and features rather than bounded, confidence-weighted information.\par}

{\small At the same time, work on safety reasoning shows that scene graphs are a useful input for detecting hazardous situations. SafetyDetect~\cite{mullen2024safetydetect}, for example, uses object relations from a scene graph to let a language model classify situations as normal or dangerous. The authors note that their approach is limited mainly by how the scene graph is created, since risk assessment can only consider what the scene graph actually contains. The quality of the underlying scene graph is therefore a key factor for any risk assessment built on top of it.\par}

{\small The existing risk-annotated scene graph framework~\cite{previous_rsg_thesis} brings these two directions together by adding risk-related information to a scene graph generated from RGB-D data. It is, however, limited to local, frame-based scene understanding and is exposed through a service interface that is not designed for robotic middleware. It also has no mechanism for keeping objects consistent across a larger environment, and it provides only basic handling of perception uncertainty. There is therefore a need for a risk-annotated scene graph that combines the spatial consistency of systems like Hydra with flexible object identification and risk-relevant object information, and that is designed for integration into a mobile robotic system during operation.\par}

\section{Goal of the Thesis}

{\small The goal of this thesis is to extend the existing risk-annotated scene graph framework~\cite{previous_rsg_thesis} into a 3D spatial scene graph that provides up-to-date environmental context for dynamic risk assessment. The scene graph should describe the objects in the environment, their spatial context, and their risk-relevant properties. The work focuses on four aspects: portability, spatial scalability, uncertainty awareness, and near-real-time operation.\par}

{\small The first goal is to integrate the baseline framework into ROS~2~\cite{macenski2022ros2}, so that it can interact directly with robotic sensors, localisation, and mapping components. The second goal is to ground the scene graph in a spatial map, so that objects observed at different times and locations share a common reference frame and remain spatially consistent when areas are revisited. Object identification should support open-vocabulary semantic perception, removing the dependence on ground-truth semantic labels or a fixed set of object classes.\par}

{\small The third goal is to represent uncertainty associated with object identification, mobility, and continued presence in the environment. The fourth goal is to support near-real-time operation on embedded robotic hardware. Finally, the framework is evaluated on the simulated uHumans2 dataset~\cite{rosinol2021kimera} and the real-world OpenLORIS-Scene dataset~\cite{shi2020openloris}, focusing on object tracking and labelling, scene graph consistency, and runtime behaviour.\par}

\section{Contributions}

{\small This thesis makes the following contributions:\par}

\begin{enumerate}
    \item {\small \textbf{A ROS~2-based risk-annotated scene graph framework.} The baseline framework is adapted for ROS~2, enabling direct integration with the sensors and software components of a robotic system. This provides a foundation for using risk-relevant scene representations within a broader robotic software architecture.\par}

    \item {\small \textbf{A spatial scene graph with open-vocabulary object identification.} The framework extends Hydra~\cite{hughes2022hydra} with flexible object identification beyond fixed semantic classes. The resulting 3D scene graph combines spatial, semantic, and mobility information while maintaining a consistent representation of the environment over time.\par}

    \item {\small \textbf{A persistent object representation with uncertainty information.} Objects are represented together with confidence information for their semantic identity, mobility, and continued presence in the environment. This allows the scene graph to express both what is currently believed about an object and the uncertainty associated with that information.\par}

    \item {\small \textbf{A design for near-real-time operation on embedded robotic hardware.} The framework is designed to maintain spatial scene understanding while incorporating computationally demanding semantic information. This allows the scene representation to remain suitable for applications in which environmental information must be updated continuously during robot operation.\par}

    \item {\small \textbf{An experimental evaluation on simulated and real-world data.} The extended framework is evaluated using the simulated uHumans2 dataset~\cite{rosinol2021kimera} and the real-world OpenLORIS-Scene dataset~\cite{shi2020openloris}. The evaluation examines object tracking and labelling, spatial scene graph consistency, and runtime behaviour to identify both the capabilities and limitations of the proposed approach.\par}
\end{enumerate}
